\chapter{2013年辽宁卷（理）}

\section{单选题}

\begin{problem}
复数 \(z = \frac{1}{\i - 1}\) 的模为（\quad）

\choices
    {\(\frac{1}{2}\)}
    {\(\frac{\sqrt{2}}{2}\)}
    {\(\sqrt{2}\)}
    {\(2\)}
\end{problem}

\begin{answer}
B
\end{answer}

\begin{solution}
由复数模的运算性质，
\(\lvert z\rvert=\dfrac{1}{\lvert\i-1\rvert}=\dfrac{1}{\sqrt{(-1)^2+1^2}}=\dfrac{\sqrt{2}}{2}\)，故选 B.
\end{solution}

\begin{problem}
已知集合 \(A = \{x \mid 0 < \log_4 x < 1\}\)，\(B = \{x \mid x \leqslant 2\}\)，则 \(A \cap B =\)（\quad）

\choices
    {\((0,1)\)}
    {\((0,2]\)}
    {\((1,2)\)}
    {\((1,2]\)}
\end{problem}

\begin{answer}
D
\end{answer}

\begin{solution}
因为底数 \(4>1\)，所以 \(0<\log_4x<1\) 等价于 \(1<x<4\)，从而 \(A=(1,4)\). 又 \(B=(-\infty,2]\)，故 \(A\cap B=(1,2]\)，选 D.
\end{solution}

\begin{problem}
已知点 \(A(1,3)\)，\(B(4,-1)\)，则与向量 \(\overrightarrow{AB}\) 同方向的单位向量为（\quad）

\choices
    {\(\left(\frac{3}{5}, -\frac{4}{5}\right)\)}
    {\(\left(\frac{4}{5}, -\frac{3}{5}\right)\)}
    {\(\left(-\frac{3}{5}, \frac{4}{5}\right)\)}
    {\(\left(-\frac{4}{5}, \frac{3}{5}\right)\)}
\end{problem}

\begin{answer}
A
\end{answer}

\begin{solution}
\(\overrightarrow{AB}=(4-1,-1-3)=(3,-4)\)，且 \(\lvert\overrightarrow{AB}\rvert=\sqrt{3^2+(-4)^2}=5\). 因此与它同方向的单位向量为
\(\dfrac{\overrightarrow{AB}}{\lvert\overrightarrow{AB}\rvert}=\left(\dfrac35,-\dfrac45\right)\)，故选 A.
\end{solution}

\begin{problem}
下面是关于公差 \(d > 0\) 的等差数列 \(\{a_n\}\) 的四个命题：

\begin{circlelist}
\item 数列 \(\{a_n\}\) 是递增数列
\item 数列 \(\{na_n\}\) 是递增数列
\item 数列 \(\left\{\frac{a_n}{n}\right\}\) 是递增数列
\item 数列 \(\{a_n + 3nd\}\) 是递增数列
\end{circlelist}

其中的真命题为（\quad）

\choices
    {\(p_1\)，\(p_2\)}
    {\(p_3\)，\(p_4\)}
    {\(p_2\)，\(p_3\)}
    {\(p_1\)，\(p_4\)}
\end{problem}

\begin{answer}
D
\end{answer}

\begin{solution}
设 \(a_n=a_1+(n-1)d\)，其中 \(d>0\). 对命题 \(p_1\)，有 \(a_{n+1}-a_n=d>0\)，所以 \(p_1\) 为真.

对 \(p_2\)，有
\((n+1)a_{n+1}-na_n=a_n+(n+1)d=a_1+2nd\)，它不一定为正，例如取 \(a_1=-10\)，\(d=1\)，\(n=1\) 时为 \(-8\)，所以 \(p_2\) 不一定为真.

对 \(p_3\)，有
\(\dfrac{a_{n+1}}{n+1}-\dfrac{a_n}{n}=\dfrac{d-a_1}{n(n+1)}\)，它不一定为正，例如取 \(a_1=2\)，\(d=1\)，所以 \(p_3\) 不一定为真.

对 \(p_4\)，有 \((a_{n+1}+3(n+1)d)-(a_n+3nd)=4d>0\)，所以 \(p_4\) 为真. 真命题为 \(p_1\)，\(p_4\)，故选 D.
\end{solution}

\begin{problem}
某班的全体学生参加英语测试，成绩的频率分布直方图如图，数据的分组依次为：\([20,40)\)，\([40,60)\)，\([60,80)\)，\([80,100]\). 若低于 \(60\) 分的人数是 \(15\)，则该班的学生人数是（\quad）

\texfigure[max width=0.46\linewidth]{img_repaint/2013/liaoning_science/q5_fig1.tex}

\choices
    {\(45\)}
    {\(50\)}
    {\(55\)}
    {\(60\)}
\end{problem}

\begin{answer}
B
\end{answer}

\begin{solution}
由直方图可读出 \([20,40)\) 与 \([40,60)\) 两组的频率密度分别为 \(0.005\) 与 \(0.010\)，两组组距均为 \(20\)，所以低于 \(60\) 分的频率为
\(20\times0.005+20\times0.010=0.3\).

设全班人数为 \(N\)，则 \(0.3N=15\)，解得 \(N=50\)，故选 B.
\end{solution}

\begin{problem}
在 \(\triangle ABC\) 中，内角 \(A\)，\(B\)，\(C\) 所对的边分别为 \(a\)，\(b\)，\(c\)，若 \(a \sin B \cos C + c \sin B \cos A = \frac{1}{2} b\)，且 \(a > b\)，则 \(\angle B =\)（\quad）

\choices
    {\(\frac{\pi}{6}\)}
    {\(\frac{\pi}{3}\)}
    {\(\frac{2\pi}{3}\)}
    {\(\frac{5\pi}{6}\)}
\end{problem}

\begin{answer}
A
\end{answer}

\begin{solution}
由正弦定理，设外接圆半径为 \(R\)，则 \(a=2R\sin A\)，\(b=2R\sin B\)，\(c=2R\sin C\). 因而
\[
a\cos C+c\cos A=2R(\sin A\cos C+\sin C\cos A)=2R\sin(A+C)=2R\sin B=b
\]
题设等式化为 \(b\sin B=\dfrac12b\)，所以 \(\sin B=\dfrac12\). 于是 \(B=\dfrac\pi6\) 或 \(B=\dfrac{5\pi}6\).

若 \(B=\dfrac{5\pi}6\)，则 \(A+C=\dfrac\pi6\)，从而 \(A<B\)，由正弦定理得 \(a<b\)，与题设矛盾. 所以 \(B=\dfrac\pi6\)，故选 A.
\end{solution}

\begin{problem}
使 \(\left(3x + \frac{1}{x\sqrt{x}}\right)^n\)（\(n \in \mathbf{N}_{+}\)）的展开式中含有常数项的最小的 \(n\) 为（\quad）

\choices
    {\(4\)}
    {\(5\)}
    {\(6\)}
    {\(7\)}
\end{problem}

\begin{answer}
B
\end{answer}

\begin{solution}
二项式通项为
\(\mathrm C_n^k(3x)^{n-k}\left(x^{-3/2}\right)^k\)，其中 \(0\leqslant k\leqslant n\). 该项中 \(x\) 的指数为
\(n-k-\dfrac32k=n-\dfrac52k\).

含常数项需满足 \(n-\dfrac52k=0\)，即 \(2n=5k\). 因而 \(n\) 必须是 \(5\) 的倍数，最小正整数取 \(n=5\)，此时 \(k=2\) 符合 \(0\leqslant k\leqslant n\)，故选 B.
\end{solution}

\begin{problem}
执行如图所示的程序框图，若输入 \(n = 10\)，则输出 \(S =\)（\quad）

\texfigure[max width=0.26\linewidth]{img_repaint/2013/liaoning_science/q8_fig1.tex}

\choices
    {\(\frac{5}{11}\)}
    {\(\frac{10}{11}\)}
    {\(\frac{36}{55}\)}
    {\(\frac{72}{55}\)}
\end{problem}

\begin{answer}
A
\end{answer}

\begin{solution}
程序中 \(i\) 依次取 \(2\)，\(4\)，\(6\)，\(8\)，\(10\)，所以
\[
S=\frac1{2^2-1}+\frac1{4^2-1}+\frac1{6^2-1}+\frac1{8^2-1}+\frac1{10^2-1}
\]
利用 \(\dfrac1{i^2-1}=\dfrac12\left(\dfrac1{i-1}-\dfrac1{i+1}\right)\)，得
\(S=\dfrac12\left(1-\dfrac1{11}\right)=\dfrac5{11}\)，故选 A.
\end{solution}

\begin{problem}
已知点 \(O(0,0)\)，\(A(0,b)\)，\(B(a,a^{3})\). 若 \(\triangle OAB\) 为直角三角形，则必有（\quad）

\choices
    {\(b = a^{3}\)}
    {\(b = a^{3} + \frac{1}{a}\)}
    {\(\left(b - a^{3}\right)\left(b - a^{3} - \frac{1}{a}\right) = 0\)}
    {\(\left|b - a^3\right| + \left|b - a^3 - \frac{1}{a}\right| = 0\)}
\end{problem}

\begin{answer}
C
\end{answer}

\begin{solution}
三点构成非退化三角形，故 \(a\ne0\)，\(b\ne0\). 直角若在 \(O\)，则
\(\overrightarrow{OA}\cdot\overrightarrow{OB}=ba^3\ne0\)，不可能.

若直角在 \(A\)，则
\[
\overrightarrow{AO}\cdot\overrightarrow{AB}=(0,-b)\cdot(a,a^3-b)=b(b-a^3)=0
\]
所以 \(b=a^3\).

若直角在 \(B\)，则
\[
\overrightarrow{BO}\cdot\overrightarrow{BA}=(-a,-a^3)\cdot(-a,b-a^3)
=a^2-a^3(b-a^3)=a^2\bigl(1-ab+a^4\bigr)=0
\]
所以 \(b=a^3+\dfrac1a\). 因而必有
\(\left(b-a^3\right)\left(b-a^3-\dfrac1a\right)=0\)，故选 C.
\end{solution}

\begin{problem}
已知直三棱柱 \(ABC-A_1B_1C_1\) 的 \(6\) 个顶点都在球 \(O\) 的球面上. 若 \(AB = 3\)，\(AC = 4\)，\(AB \perp AC\)，\(AA_1 = 12\)，则球 \(O\) 的半径为（\quad）

\choices
    {\(\frac{3\sqrt{17}}{2}\)}
    {\(2\sqrt{10}\)}
    {\(\frac{13}{2}\)}
    {\(3\sqrt{10}\)}
\end{problem}

\begin{answer}
C
\end{answer}

\begin{solution}
取坐标 \(A=(0,0,0)\)，\(B=(3,0,0)\)，\(C=(0,4,0)\)，\(A_1=(0,0,12)\). 设球心为 \(O=(u,v,w)\). 由 \(OA=OB\)、\(OA=OC\)、\(OA=OA_1\) 分别得
\(u=\dfrac32\)，\(v=2\)，\(w=6\).

因此球半径满足
\[
R^2=OA^2=u^2+v^2+w^2=\left(\frac32\right)^2+2^2+6^2=\frac{169}{4}
\]
所以 \(R=\dfrac{13}{2}\)，故选 C.
\end{solution}

\begin{problem}
已知函数 \(f(x) = x^2 - 2(a + 2)x + a^2\)，\(g(x) = -x^2 + 2(a - 2)x - a^2 + 8\)，设 \(H_1(x) = \max\{f(x), g(x)\}\)，\(H_2(x) = \min\{f(x), g(x)\}\)（\(\max\{p, q\}\) 表示 \(p\)，\(q\) 中的较大值，\(\min\{p, q\}\) 表示 \(p\)，\(q\) 中的较小值）. 记 \(H_1(x)\) 的最小值为 \(A\)，\(H_2(x)\) 的最大值为 \(B\)，则 \(A - B =\)（\quad）

\choices
    {\(16\)}
    {\(-16\)}
    {\(a^2 - 2a - 16\)}
    {\(a^2 + 2a - 16\)}
\end{problem}

\begin{answer}
B
\end{answer}

\begin{solution}
两函数之差为
\[
f(x)-g(x)=2\bigl((x-a)^2-4\bigr)
\]
所以当 \(\lvert x-a\rvert<2\) 时 \(f(x)<g(x)\)，当 \(\lvert x-a\rvert>2\) 时 \(f(x)>g(x)\)，交点为 \(x=a-2\)，\(a+2\). 又 \(f(x)=(x-a-2)^2-4a-4\)，\(g(x)=-(x-a+2)^2-4a+12\).

在区间 \([a-2,a+2]\) 内，\(H_1=g\)，其最小值在端点取得，比较端点值 \(12-4a\) 与 \(-4a-4\)，得 \(A=-4a-4\). 在区间外 \(H_1=f\)，其值不小于相应端点值，故不改变这个最小值.

在区间 \([a-2,a+2]\) 内，\(H_2=f\)，其最大值为端点较大值 \(12-4a\). 区间外 \(H_2=g\)，且其最大值也不超过端点值，因此 \(B=12-4a\). 所以 \(A-B=-16\)，故选 B.
\end{solution}

\begin{problem}
设函数 \( f(x) \) 满足 \(x^{2}f'(x) + 2xf(x) = \frac{\e^{x}}{x}\)，\(f(2) = \frac{\e^{2}}{8}\)，则 \( x > 0 \) 时，\( f(x) \)（\quad）

\choices
    {有极大值，无极小值}
    {有极小值，无极大值}
    {既有极大值又有极小值}
    {既无极大值也无极小值}
\end{problem}

\begin{answer}
D
\end{answer}

\begin{solution}
令 \(H(x)=x^2f(x)\)，则由题设
\(H'(x)=x^2f'(x)+2xf(x)=\dfrac{\e^x}{x}\). 再令
\(G(x)=\e^x-2H(x)\)，则
\[
G'(x)=\e^x-2\frac{\e^x}{x}=\frac{\e^x(x-2)}{x}
\]
所以 \(G\) 在 \((0,2)\) 上递减，在 \((2,+\infty)\) 上递增. 由 \(f(2)=\dfrac{\e^2}{8}\) 得 \(H(2)=\dfrac{\e^2}{2}\)，故 \(G(2)=0\)，从而 \(G(x)\geqslant0\)，且除 \(x=2\) 外均严格为正.

将原方程乘以 \(x\)，得 \(x^3f'(x)=\e^x-2x^2f(x)=G(x)\). 因为 \(x>0\)，所以 \(f'(x)\geqslant0\)，且函数在任意非退化区间上严格递增；在 \(x=2\) 处导数为零但不改变单调性. 因此 \(f(x)\) 既无极大值也无极小值，故选 D.
\end{solution}

\section{填空题}

\begin{problem}
某几何体的三视图如图所示，则该几何体的体积是\fillinblank{}.

\bitmapfigure[max width=0.46\linewidth]{img/2013/liaoning_science/q13_fig1.png}
\end{problem}

\begin{answer}
\(16\pi - 16\)
\end{answer}

\begin{solution}
由三视图可知，该几何体是一个底面半径为 \(2\)、高为 \(4\) 的圆柱，内部挖去一个底面为边长 \(2\) 的正方形、高为 \(4\) 的直棱柱. 因而体积为
\(V=\pi\times2^2\times4-2^2\times4=16\pi-16\).
\end{solution}

\begin{problem}
已知等比数列 \(\{a_n\}\) 是递增数列，\(S_n\) 是 \(\{a_n\}\) 的前 \(n\) 项和，若 \(a_1\)，\(a_3\) 是方程 \(x^2 - 5x + 4 = 0\) 的两个根，则 \(S_6 =\)\fillinblank{}.
\end{problem}

\begin{answer}
\(63\)
\end{answer}

\begin{solution}
方程 \(x^2-5x+4=0\) 的两个根为 \(1\)，\(4\). 设等比数列公比为 \(q\)，则 \(a_3=a_1q^2\).

若 \(a_1=1\)，\(a_3=4\)，则 \(q^2=4\)，所以 \(q=2\) 或 \(q=-2\). 当 \(q=2\) 时数列为 \(1\)，\(2\)，\(4\)，\(\ldots\)，递增；当 \(q=-2\) 时 \(a_2=-2<a_1\)，不递增. 因而此时只能取 \(q=2\).

若 \(a_1=4\)，\(a_3=1\)，则 \(q^2=\dfrac14\). 取 \(q=\dfrac12\) 时数列递减，取 \(q=-\dfrac12\) 时各项交替变号，均不可能递增. 因而 \(a_1=1\)，\(q=2\)，所以
\(S_6=1+2+2^2+\cdots+2^5=\dfrac{2^6-1}{2-1}=63\).
\end{solution}

\begin{problem}
已知椭圆 \(C:\frac{x^2}{a^2} + \frac{y^2}{b^2} = 1\)（\(a > b > 0\)）的左焦点为 \(F\)，\(C\) 与过原点的直线相交于 \(A\)，\(B\) 两点，连接 \(AF\)，\(BF\)，若 \(|AB| = 10\)，\(|AF| = 6\)，\(\cos \angle ABF = \frac{4}{5}\)，则椭圆 \(C\) 的离心率 \(e =\)\fillinblank{}.
\end{problem}

\begin{answer}
\(\frac{5}{7}\)
\end{answer}

\begin{solution}
设椭圆的右焦点为 \(F'\)，因为直线 \(AB\) 过原点，故原点是弦 \(AB\) 的中点，且 \(AB=10\). 由余弦定理，在三角形 \(ABF\) 中设 \(BF=t\)，有
\[
6^2=10^2+t^2-2\times10\times t\times\frac45
\]
即 \((t-8)^2=0\)，所以 \(BF=8\). 此时三角形 \(ABF\) 的三边为 \(6\)，\(8\)，\(10\)，且 \(\angle AFB=90^\circ\). 原点为斜边 \(AB\) 的中点，所以 \(OF=OA=5\)，从而焦距 \(c=5\).

又由中心对称性，\(AF'=BF=8\). 根据椭圆定义，\(AF+AF'=2a\)，得 \(2a=6+8=14\)，即 \(a=7\). 因而离心率
\(e=\dfrac ca=\dfrac57\).
\end{solution}

\begin{problem}
为了考察某校各班参加课外书法小组的人数，从全校随机抽取 \(5\) 个班级，把每个班级参加该小组的人数作为样本数据. 已知样本平均数为 \(7\)，样本方差为 \(4\)，且样本数据互不相同，则样本数据中的最大值为\fillinblank{}.
\end{problem}

\begin{answer}
\(10\)
\end{answer}

\begin{solution}
设样本数据为 \(x_1\)，\(x_2\)，\(x_3\)，\(x_4\)，\(x_5\)，则
\[
\sum_{i=1}^5x_i=5\times7=35
\]
\[
\sum_{i=1}^5(x_i-7)^2=5\times4=20
\]
所以 \(\sum x_i^2=20+5\times7^2=265\). 设最大值为 \(m\)，则其余四个数据的和与平方和分别为 \(35-m\) 和 \(265-m^2\). 由均方不小于平均数的平方，
\[
265-m^2\geqslant\frac{(35-m)^2}{4}
\]
整理得 \((m-3)(m-11)\leqslant0\)，所以 \(m\leqslant11\). 若 \(m=11\)，上式取等号，说明其余四个数据都等于 \(\dfrac{35-11}{4}=6\)，与样本数据互不相同矛盾，故 \(m\leqslant10\).

数据 \(4\)，\(6\)，\(7\)，\(8\)，\(10\) 的和为 \(35\)，平方和为 \(265\)，满足题设条件，故最大值可以取到 \(10\).
\end{solution}

\section{解答题}

\begin{problem}
设向量 \(\bs{a} = (\sqrt{3}\sin x, \sin x)\)，\(\bs{b} = (\cos x, \sin x)\)，\(x \in \left[0, \frac{\pi}{2}\right]\).

\begin{enumerate}
\item 若 \(\left|\bs{a}\right| = \left|\bs{b}\right|\)，求 \(x\) 的值；
\item 设函数 \(f(x) = \bs{a} \cdot \bs{b}\)，求 \(f(x)\) 的最大值.
\end{enumerate}
\end{problem}

\begin{solution}
\begin{enumerate}
\item 由向量模相等得
\[
\left|\bs{a}\right|^2=3\sin^2x+\sin^2x=4\sin^2x
\]
\[
\left|\bs{b}\right|^2=\cos^2x+\sin^2x=1
\]
所以 \(\sin^2x=\dfrac14\). 因为 \(x\in\left[0,\dfrac\pi2\right]\)，故 \(\sin x=\dfrac12\)，从而 \(x=\dfrac\pi6\).

\item 由数量积的坐标公式，
\[
f(x)=\bs{a}\cdot\bs{b}=\sqrt3\sin x\cos x+\sin^2x
=\frac{\sqrt3}{2}\sin2x+\frac{1-\cos2x}{2}
=\sin\left(2x-\frac\pi6\right)+\frac12
\]
当 \(x\in\left[0,\dfrac\pi2\right]\) 时，\(2x-\dfrac\pi6\) 的取值区间为 \(\left[-\dfrac\pi6,\dfrac{5\pi}6\right]\)，其中包含使正弦值为 \(1\) 的角 \(\dfrac\pi2\)，对应 \(x=\dfrac\pi3\). 因此 \(f(x)\) 的最大值为 \(\dfrac32\).
\end{enumerate}
\end{solution}

\begin{problem}
如图，\(AB\) 是圆的直径，\(PA\) 垂直圆所在的平面，\(C\) 是圆上的点.

\begin{enumerate}
\item 求证：平面 \(PAC \perp PBC\)；
\item 若 \(AB = 2\)，\(AC = 1\)，\(PA = 1\)，求二面角 \(C - PB - A\) 的余弦值.
\end{enumerate}

\FigureLayoutDeclare{img/2013/liaoning_science/q18_fig1.png}{8.0cm}{72.966bp 110.41bp 12.481bp 12.001bp}
\bitmapfigure[max width=0.38\linewidth]{img/2013/liaoning_science/q18_fig1.png}
\end{problem}

\begin{solution}
\begin{enumerate}
\item 因为 \(AB\) 是圆的直径，且 \(C\) 在圆上，所以 \(\angle ACB=90^\circ\)，即 \(AC\perp BC\). 又因为 \(PA\) 垂直于圆所在平面，故 \(PA\perp BC\). 直线 \(AC\)，\(PA\) 是平面 \(PAC\) 内相交的两条直线，且 \(BC\) 同时垂直于它们，所以 \(BC\perp\) 平面 \(PAC\). 由于 \(BC\subset\) 平面 \(PBC\)，故平面 \(PAC\perp\) 平面 \(PBC\).

\item 在直角三角形 \(ABC\) 中，\(BC=\sqrt{AB^2-AC^2}=\sqrt3\). 建立空间直角坐标系，取
\(A=(0,0,0)\)，\(B=(2,0,0)\)，\(C=\left(\dfrac12,\dfrac{\sqrt3}{2},0\right)\)，\(P=(0,0,1)\).

平面 \(PBA\) 的一个法向量为 \(\bs{n}_1=(0,1,0)\). 平面 \(PBC\) 的两个方向向量可取 \(\overrightarrow{PB}=(2,0,-1)\)，\(\overrightarrow{PC}=\left(\dfrac12,\dfrac{\sqrt3}{2},-1\right)\)，故其法向量可取
\[
\bs{n}_2=\overrightarrow{PB}\times\overrightarrow{PC}=\left(\frac{\sqrt3}{2},\frac32,\sqrt3\right)
\]
两平面所成锐二面角的余弦为
\[
\frac{\lvert\bs{n}_1\cdot\bs{n}_2\rvert}{\lvert\bs{n}_1\rvert\lvert\bs{n}_2\rvert}
=\frac{\frac32}{\sqrt{\frac34+\frac94+3}}=\frac{\sqrt6}{4}
\]
\end{enumerate}
\end{solution}

\begin{problem}
现有 \(10\) 道题，其中 \(6\) 道甲类题，\(4\) 道乙类题，张同学从中任取 \(3\) 道题解答.

\begin{enumerate}
\item 求张同学至少取到 \(1\) 道乙类题的概率；
\item 已知所取的 \(3\) 道题中有 \(2\) 道甲类题，\(1\) 道乙类题. 设张同学答对每道甲类题的概率都是 \(\frac{3}{5}\)，答对每道乙类题的概率都是 \(\frac{4}{5}\)，且各题答对与否相互独立. 用 \(X\) 表示张同学答对题的个数，求 \(X\) 的分布列和数学期望.
\end{enumerate}
\end{problem}

\begin{solution}
\begin{enumerate}
\item 从 \(10\) 道题中任取 \(3\) 道共有 \(\mathrm C_{10}^3\) 种取法. 至少取到一道乙类题的对立事件是三道全为甲类题，故所求概率为
\[
1-\frac{\mathrm C_6^3}{\mathrm C_{10}^3}=1-\frac{20}{120}=\frac56
\]

\item 设两道甲类题答对的个数为 \(Y\)，则 \(Y\) 服从参数为 \(2\)，\(\dfrac35\) 的二项分布；乙类题答对的概率为 \(\dfrac45\). 随机变量 \(X\) 的可能取值为 \(0\)，\(1\)，\(2\)，\(3\)，分别计算得
\[
P(X=0)=\left(\frac25\right)^2\frac15=\frac4{125}
\]
\[
P(X=1)=2\frac35\frac25\frac15+\left(\frac25\right)^2\frac45=\frac{28}{125}
\]
\[
P(X=2)=\left(\frac35\right)^2\frac15+2\frac35\frac25\frac45=\frac{57}{125}
\]
\[
P(X=3)=\left(\frac35\right)^2\frac45=\frac{36}{125}
\]
因此分布列为

\begin{center}
\begin{tabular}{c|cccc}
\(X\) & \(0\) & \(1\) & \(2\) & \(3\) \\
\hline
\(P\) & \(\dfrac4{125}\) & \(\dfrac{28}{125}\) & \(\dfrac{57}{125}\) & \(\dfrac{36}{125}\)
\end{tabular}
\end{center}

数学期望为
\[
E(X)=0\cdot\frac4{125}+1\cdot\frac{28}{125}+2\cdot\frac{57}{125}+3\cdot\frac{36}{125}=2
\]
\end{enumerate}
\end{solution}

\begin{problem}
如图，抛物线 \(C_1:x^2 = 4y\)，\(C_2:x^2 = -2py\)（\(p > 0\)）. 点 \(M(x_0, y_0)\) 在抛物线 \(C_2\) 上，过 \(M\) 作 \(C_1\) 的切线，切点为 \(A\)，\(B\)（\(M\) 为原点 \(O\) 时，\(A\)，\(B\) 重合于 \(O\)）. 当 \(x_0 = 1 - \sqrt{2}\) 时，切线 \(MA\) 的斜率为 \(-\frac{1}{2}\).

\begin{enumerate}
\item 求 \(p\) 的值；
\item 当 \(M\) 在 \(C_2\) 上运动时，求线段 \(AB\) 中点 \(N\) 的轨迹方程（\(A\)，\(B\) 重合于 \(O\) 时，中点为 \(O\)）.
\end{enumerate}

\bitmapfigure[max width=0.34\linewidth]{img/2013/liaoning_science/q20_fig1.png}
\end{problem}

\begin{solution}
\begin{enumerate}
\item 将 \(C_1\) 写成 \(y=\dfrac{x^2}{4}\). 设 \(C_1\) 上切点对应的切线斜率为 \(m\)，则切点为 \((2m,m^2)\)，切线方程为 \(y=mx-m^2\). 过 \(M(x_0,y_0)\) 的切线斜率满足
\[
y_0=mx_0-m^2,\qquad\text{即}\qquad m^2-x_0m+y_0=0
\]
题设一条切线斜率为 \(-\dfrac12\)，所以 \(y_0=-\dfrac12x_0-\dfrac14\). 当 \(x_0=1-\sqrt2\) 时，\(y_0=-\dfrac{x_0^2}{4}\). 又 \(M\in C_2\)，所以 \(y_0=-\dfrac{x_0^2}{2p}\)，故 \(p=2\).

\item 由上问，\(C_2\) 的方程为 \(y=-\dfrac{x^2}{4}\). 设两条切线的斜率为 \(m_1\)，\(m_2\)，则由切线斜率方程和韦达定理，
\[
m_1+m_2=x_0
\]
\[
m_1m_2=y_0=-\frac{x_0^2}{4}
\]
两切点分别为 \(A=(2m_1,m_1^2)\)，\(B=(2m_2,m_2^2)\)，故中点 \(N=(x_N,y_N)\) 满足
\[
x_N=m_1+m_2=x_0
\]
\[
y_N=\frac{m_1^2+m_2^2}{2}=\frac{(m_1+m_2)^2-2m_1m_2}{2}=\frac{3x_0^2}{4}
\]
消去 \(x_0\) 得 \(x_N^2=\dfrac43y_N\). 当 \(M=O\) 时 \(x_0=0\)，中点也是 \(O\)，同样满足该方程. 因而轨迹方程为 \(x^2=\dfrac43y\).
\end{enumerate}
\end{solution}

\begin{problem}
已知函数 \(f(x) = (1 + x)\e^{-2x}\)，\(g(x) = ax + \frac{x^3}{2} + 1 + 2x\cos x\)，当 \(x \in [0,1]\) 时，

\begin{enumerate}
\item 求证：\(1 - x \leqslant f(x) \leqslant \frac{1}{1 + x}\)；
\item 若 \(f(x) \geqslant g(x)\) 恒成立，求实数 \(a\) 的取值范围.
\end{enumerate}
\end{problem}

\begin{solution}
\begin{enumerate}
\item 令 \(\phi(x)=(1+x)\e^{-2x}-(1-x)\). 则 \(\phi(0)=0\)，且 \(\phi'(x)=1-(1+2x)\e^{-2x}\). 令 \(u(x)=(1+2x)\e^{-2x}\)，则 \(u'(x)=-4x\e^{-2x}\leqslant0\)，所以在 \([0,1]\) 上 \(u(x)\leqslant u(0)=1\)，从而 \(\phi'(x)\geqslant0\). 因此 \(\phi(x)\geqslant0\)，即 \(f(x)\geqslant1-x\).

又因为 \(\e^x\geqslant1+x>0\)，所以 \(\e^{-2x}\leqslant\dfrac1{(1+x)^2}\)，从而 \(f(x)=(1+x)\e^{-2x}\leqslant\dfrac1{1+x}\). 两个不等式均得证.

\item 当 \(x>0\) 时，\(f(x)\geqslant g(x)\) 等价于
\[
a\leqslant\Phi(x):=\frac{(1+x)\e^{-2x}-1-\frac{x^3}{2}-2x\cos x}{x}
\]
若不等式恒成立，则对所有 \(x>0\) 有 \(a\leqslant\Phi(x)\). 由洛必达法则，
\[
\lim_{x\to0^+}\Phi(x)=\left[(1+x)\e^{-2x}\right]'_{x=0}-2\cos0=-1-2=-3
\]
故必有 \(a\leqslant-3\).

反之，设 \(a\leqslant-3\). 因为 \(x\geqslant0\)，有 \(g(x)\leqslant g_{-3}(x):=-3x+\dfrac{x^3}{2}+1+2x\cos x\). 由第（1）问，只需证明 \(1-x\geqslant g_{-3}(x)\). 而
\[
1-x-g_{-3}(x)=x\left(2(1-\cos x)-\frac{x^2}{2}\right)
\]
令 \(r(x)=1-\cos x-\dfrac{x^2}{4}\). 由 \(1-\cos x=2\sin^2\dfrac{x}{2}\leqslant\dfrac{x^2}{2}\)，可得 \(\cos x\geqslant1-\dfrac{x^2}{2}\geqslant\dfrac12\)（这里 \(0\leqslant x\leqslant1\)），所以 \(r''(x)=\cos x-\dfrac12\geqslant0\). 再由 \(r(0)=r'(0)=0\)，得 \(r(x)\geqslant0\). 因而 \(1-x-g_{-3}(x)=2xr(x)\geqslant0\)，于是 \(f(x)\geqslant g(x)\) 恒成立.

所以实数 \(a\) 的取值范围为 \((-\infty,-3]\).
\end{enumerate}
\end{solution}

\section{选考题：解答题}

\begin{problem}
如图，\(AB\) 为 \(\odot O\) 直径，直线 \(CD\) 与 \(\odot O\) 相切于 \(E\)，\(AD\) 垂直 \(CD\) 于 \(D\)，\(BC\) 垂直 \(CD\) 于 \(C\)，\(EF\) 垂直 \(AB\) 于 \(F\)，连接 \(AE\)，\(BE\). 证明：

\begin{enumerate}
\item \(\angle FEB = \angle CEB\)；
\item \(EF^2 = AD \cdot BC\).
\end{enumerate}

\bitmapfigure[max width=0.43\linewidth]{img/2013/liaoning_science/q22_fig1.png}
\end{problem}

\begin{solution}
\begin{enumerate}
\item 因为 \(AB\) 是直径，所以 \(\angle AEB=90^\circ\). 由弦切角定理，\(\angle CEB=\angle EAB\). 又 \(EF\perp AB\)，所以 \(\angle FEB=90^\circ-\angle EBA\). 在直角三角形 \(AEB\) 中，\(\angle EAB+\angle EBA=90^\circ\)，故 \(\angle FEB=\angle EAB=\angle CEB\).

\item 设 \(\alpha=\angle EAB\)，\(\beta=\angle EBA\)，则 \(\alpha+\beta=90^\circ\). 由弦切角定理，\(\angle AED=\beta\)，\(\angle BEC=\alpha\). 在两个直角三角形 \(AED\)，\(BCE\) 中，
\[
AD=AE\sin\beta
\]
\[
BC=BE\sin\alpha
\]
又在直角三角形 \(AEB\) 中，\(AE=AB\cos\alpha\)，\(BE=AB\sin\alpha\)，且从 \(E\) 向斜边 \(AB\) 作的高满足 \(EF=AB\sin\alpha\cos\alpha\). 因此
\[
AD=AB\cos\alpha\sin\beta=AB\cos^2\alpha=EF\cot\alpha
\]
\[
BC=AB\sin\alpha\sin\alpha=EF\tan\alpha
\]
于是 \(AD\cdot BC=EF^2\cot\alpha\tan\alpha=EF^2\).
\end{enumerate}
\end{solution}

\begin{problem}
在直角坐标系 \(xOy\) 中，以 \(O\) 为极点，\(x\) 轴正半轴为极轴建立极坐标系. 圆 \(C_1\)，直线 \(C_2\) 的极坐标方程分别为 \(\rho = 4\sin \theta\)，\(\rho \cos \left(\theta - \frac{\pi}{4}\right) = 2\sqrt{2}\).

\begin{enumerate}
\item 求 \(C_1\) 与 \(C_2\) 交点的极坐标；
\item 设 \(P\) 为 \(C_1\) 的圆心，\(Q\) 为 \(C_1\) 与 \(C_2\) 交点连线的中点，已知直线 \(PQ\) 的参数方程为 \(\begin{cases}
x = t^3 + a,\\
y = \frac{b}{2}t^3 + 1,
\end{cases}\)（\(t \in \R\) 为参数），求 \(a\)，\(b\) 的值.
\end{enumerate}
\end{problem}

\begin{solution}
\begin{enumerate}
\item 由 \(x=\rho\cos\theta\)，\(y=\rho\sin\theta\)，圆 \(C_1\) 的方程化为 \(x^2+y^2=4y\). 直线 \(C_2\) 的方程化为
\(\dfrac{x+y}{\sqrt2}=2\sqrt2\)，即 \(x+y=4\).

代入 \(y=4-x\)，得
\[
x^2+(4-x)^2=4(4-x)
\]
即 \(2x(x-2)=0\). 所以交点为 \((0,4)\)，\((2,2)\)，对应的极坐标分别为 \(\left(4,\dfrac\pi2\right)\) 和 \(\left(2\sqrt2,\dfrac\pi4\right)\).

\item 圆 \(C_1\) 的普通方程为 \(x^2+(y-2)^2=4\)，故圆心 \(P=(0,2)\). 两交点的中点为 \(Q=(1,3)\)，所以直线 \(PQ\) 的方程为 \(y=x+2\).

已知参数方程满足 \(y-1=\dfrac b2(x-a)\)，即 \(y=\dfrac b2x+1-\dfrac{ab}{2}\). 与 \(y=x+2\) 比较系数和常数项，得 \(b=2\)，\(a=-1\).
\end{enumerate}
\end{solution}

\begin{problem}
已知函数 \(f(x) = |x - a|\)，其中 \(a > 1\).

\begin{enumerate}
\item 当 \(a = 2\) 时，求不等式 \(f(x) \geqslant 4 - |x - 4|\) 的解集；
\item 已知关于 \(x\) 的不等式 \(\left|f(2x + a) - 2f(x)\right| \leqslant 2\) 的解集为 \(\{x \mid 1 \leqslant x \leqslant 2\}\)，求 \(a\) 的值.
\end{enumerate}
\end{problem}

\begin{solution}
\begin{enumerate}
\item 当 \(a=2\) 时，分三段讨论.

当 \(x\leqslant2\) 时，不等式为 \(2-x\geqslant x\)，解得 \(x\leqslant1\).

当 \(2\leqslant x\leqslant4\) 时，不等式为 \(x-2\geqslant x\)，无解.

当 \(x\geqslant4\) 时，不等式为 \(x-2\geqslant8-x\)，解得 \(x\geqslant5\). 所以解集为 \((-\infty,1]\cup[5,+\infty)\).

\item 因为 \(a>1\)，有
\[
f(2x+a)=\lvert2x+a-a\rvert=2\lvert x\rvert
\]
所以题设不等式等价于 \(\bigl\lvert\lvert x\rvert-\lvert x-a\rvert\bigr\rvert\leqslant1\).

当 \(x\leqslant0\) 时，左侧等于 \(a>1\)，无解；当 \(0\leqslant x\leqslant a\) 时，不等式为 \(\lvert2x-a\rvert\leqslant1\)，即 \(\dfrac{a-1}{2}\leqslant x\leqslant\dfrac{a+1}{2}\)；当 \(x\geqslant a\) 时，左侧等于 \(a>1\)，也无解. 因而解集为 \(\left[\dfrac{a-1}{2},\dfrac{a+1}{2}\right]\).

由它等于 \([1,2]\)，得 \(\dfrac{a-1}{2}=1\)，所以 \(a=3\)，且右端点也相应为 \(2\).
\end{enumerate}
\end{solution}
